\section{Related Work}~\label{related}

In this section, we review related work from the perspective of human involvement in robot decision making. We first discuss human corrective systems and system-initiated interaction for reducing supervision burden. We then review knowledge-based robotic systems and planning under uncertainty, which provide the technical foundations for selective human assistance.

\subsection{Human Corrective Systems}~\label{rel:HCS}

Despite recent advances in robotic autonomy, robots continue to face uncertainty arising from perception errors, action failures, and environmental changes. To improve reliability, prior work has incorporated humans into the decision-making loop through shared autonomy, interactive planning, and human-on-the-loop (HOTL) systems~\cite{kelly2019hg,mandlekar2020human}. In HOTL systems, robots execute tasks autonomously while human operators monitor execution and intervene when failures occur~\cite{scharre2015introduction,nahavandi2017trusted}. This supervisory structure can improve safety and robustness because humans can compensate for failures that autonomous systems cannot resolve.

However, effective supervision requires operators to continuously monitor robot behavior and detect situations that require intervention. Situation awareness research suggests that operators must perceive relevant events, understand their implications, and determine appropriate corrective actions~\cite{endsley2018automation}. Maintaining this level of awareness over long periods can be cognitively demanding. Continuous monitoring has been associated with vigilance decrement, inattentional blindness, increased workload, and operator fatigue~\cite{davies1982psychology,davies2013human,mack2003inattentional,wong2017workload}. As a result, requiring humans to continuously identify robot failures can limit the scalability and effectiveness of long-term human--robot collaboration.

To reduce supervision burden, recent work has explored proactive robot behavior, where robots request assistance only when needed rather than relying on humans to continuously detect failures~\cite{baraglia2017efficient}. This shift transfers part of the problem detection burden from the user to the robot and motivates systems that can autonomously determine when human assistance is necessary.

\subsection{System-Initiated Querying}~\label{rel:query}

Recent research has explored uncertainty-aware mechanisms for determining when robots should request human input. For example, conformal prediction-based approaches estimate uncertainty through prediction sets and use their size to decide whether to query the user~\cite{knowno2023,liang2024introspective,mullen2025lbap}. These methods demonstrate that uncertainty estimates can reduce unnecessary human intervention while maintaining task performance.

However, existing work has primarily focused on determining when human assistance is required. Effective human--robot interaction also depends on what information should be requested from the user. Queries that provide little information may increase operator burden without improving robot performance. Therefore, system-initiated interaction should consider both the timing of human assistance and the content of the requested information.

In addition, prior studies have mainly evaluated query systems using technical metrics such as task performance, uncertainty calibration, and query efficiency. Human-centered factors such as workload, fatigue, and supervision burden have received comparatively less attention, despite their importance in collaborative robotic systems~\cite{murata2000ergonomics,borragan2017cognitive,rajavenkatanarayanan2019towards}. Consequently, system-initiated corrective interaction should be evaluated not only in terms of robot performance but also in terms of its impact on human operators.

\subsection{Knowledge-Based Reasoning under Uncertainty}~\label{rel:kb}

Knowledge-based robotic systems represent objects, attributes, and relations through symbolic structures that support interpretable reasoning and planning~\cite{brachman2004knowledge,russell1995modern}. Predicate-based representations such as STRIPS and PDDL~\cite{fikes1971strips,aeronautiques1998pddl} describe world states through explicit facts and define actions through preconditions and effects. Logical inference, commonsense reasoning, and ontology-based representations further allow robots to infer hidden information and integrate observations into a structured world model~\cite{tenorth2012knowledge,liu2020graph,kunze2014using,hasegawa2023inferring,mota2021integrated,zhang2024icorpp,daruna2019robocse}.

While these approaches provide rich semantic representations, most knowledge-based systems assume a single deterministic world state and do not explicitly represent uncertainty over alternative hypotheses~\cite{yang2018peorl,bouguerra2008monitoring}. As a result, inferred facts are often committed as definite knowledge even when observations are uncertain.

Planning under uncertainty has been extensively studied within the POMDP framework, which represents uncertainty as a belief distribution over possible states~\cite{somani2013despot,bai2015intention,cramer2021probabilistic,bravo2019use}. Since exact belief updates are computationally expensive, various approximation methods have been proposed. Particle-based approaches represent beliefs as sampled hypotheses and have been widely adopted in robotics due to their compatibility with online planning methods such as POMCP, DESPOT, and POMCPOW~\cite{silver2010monte,somani2013despot,sunberg2018pomcpow}. These methods have been successfully applied to navigation, exploration, object search, and manipulation under uncertainty~\cite{lauri2016exploration,kim2021plgrim,xiao2019objectsearch,pajarinen2022composition}.

Several studies have combined symbolic reasoning with probabilistic planning to preserve interpretability while handling uncertainty~\cite{sanner2010symbolic,zhang2015corpp,serrano2021knowledge,li2026uncertainty,tang2025tru}. However, these approaches primarily focus on autonomous decision making and do not explicitly consider how uncertainty should be communicated to humans or how human feedback should be incorporated into the decision process.

In contrast, our work combines symbolic knowledge representation, belief-based uncertainty reasoning, and system-initiated human feedback within a unified framework. The proposed system not only determines when human assistance is required but also selects what information should be requested, while evaluating both task performance and human workload.
